% =============================================================================
%  The Architecture of Outer Admissibility
%  Flagship paper --- Adequacy Architecture Project
%
%  What this paper does:
%    Tells the full story at the right altitude.
%    From motivating problem through conceptual thesis through theorem stack
%    through summit and honest limits.
%
%  What the companion note does:
%    States the solved theorem stack concisely for technically fluent readers.
%    (Adequacy_Theorem_Companion.tex)
%
%  Source of record for theorems:
%    specs/IN-PROCESS/SPEC_045_AS1/SPEC_045_AS1_A_INTEGRATED_THEOREM_STATEMENT.md
%  Full theorem-name index:
%    adequacy-architecture-lean/THEOREM_INVENTORY.md
%
%  Build:  cd paper/ && bash build_papers.sh Adequacy_Flagship_Paper.tex
%  All theorems are machine-checked in Lean 4 with 0 sorry.
% =============================================================================

\documentclass[11pt]{article}

% --- packages ----------------------------------------------------------------
\usepackage[margin=1.1in]{geometry}
\usepackage{amsmath,amsthm,amssymb}
\usepackage{mathtools}
\usepackage{hyperref}
\usepackage{microtype}
\usepackage{enumitem}
\usepackage{xcolor}
\usepackage{tcolorbox}
\tcbuselibrary{skins,breakable}
\usepackage{fancyvrb}
\usepackage{url}

% --- theorem environments ----------------------------------------------------
\theoremstyle{plain}
\newtheorem{theorem}{Theorem}[section]
\newtheorem{corollary}[theorem]{Corollary}
\newtheorem{proposition}[theorem]{Proposition}
\newtheorem{lemma}[theorem]{Lemma}

\theoremstyle{definition}
\newtheorem{definition}[theorem]{Definition}
\newtheorem{remark}[theorem]{Remark}
\newtheorem{example}[theorem]{Example}

\theoremstyle{remark}
\newtheorem{principle}{Principle}

% --- macros ------------------------------------------------------------------
\newcommand{\Gadm}{G_{\mathrm{adm}}}
\newcommand{\Upb}{U_{\mathrm{pullback}}}
\newcommand{\Ggen}{G_{\mathrm{gen}}}
\newcommand{\AbsRawSone}{\mathrm{AbsRawS1}}
\newcommand{\Stwo}{L_2}
\newcommand{\StageD}{\mathrm{StageD}}
\newcommand{\SoneLawful}{S_1^{\mathrm{lawful}}}
\newcommand{\HFinal}{H_{\mathrm{final}}}
\newcommand{\lean}[1]{\texttt{\small #1}}
\newcommand{\leanmod}[1]{\textit{\small #1}}

% --- callout box (for key claims before technical detail) --------------------
\newtcolorbox{keybox}[1][]{
  colback=blue!4, colframe=blue!35,
  arc=3pt, boxrule=0.5pt,
  fonttitle=\bfseries, #1,
  breakable
}

% --- Lean citation box -------------------------------------------------------
\newenvironment{leanbox}{%
  \VerbatimEnvironment
  \begin{tcolorbox}[colback=gray!7, colframe=gray!40,
    arc=2pt, boxrule=0.4pt, breakable]%
  \begin{Verbatim}[fontsize=\small]%
}{%
  \end{Verbatim}%
  \end{tcolorbox}%
}

% --- metadata ----------------------------------------------------------------
\title{\textbf{The Architecture of Outer Admissibility}\\[6pt]
  \large Unification, Layering, and Residual Structure\\
  in a Certificate World}
\author{Adequacy Architecture Project}
\date{2026-03-27 \quad (machine-checked Lean 4, 0 \texttt{sorry})}

% =============================================================================
\begin{document}
\maketitle

\begin{abstract}
Many formal and applied systems rely on outer certificates, summaries,
compressed witnesses, or interface-level predicates that purport to stand
in for richer internal structure.  The central structural question is:
\emph{when do these outer certificates belong to one honest admissibility
class, when are they different presentations of the same gate, when do they
unlock richer higher-layer consequences, and when do stronger local packs
remain residual rather than primitive?}

This paper answers that question within the Adequacy Architecture framework.
The main discovery is that the apparent plurality of outer admissibility
routes collapses to a single unified gate $\Gadm \equiv \Upb$, with a real
frontier consequence, a grammar-cut maximality theorem, a disciplined
boundary finding, a functorial second layer above admission, and
a residual band of stronger IC-local content that is classified rather than
dismissed.  The result is a strict three-level architecture, every boundary
of which is characterized by a named, machine-checked theorem.

The paper is structured to earn these claims.  It starts with the motivating
problem at the correct altitude, develops the conceptual thesis before the
formal machinery, and ends by honestly characterizing what lies beyond the
current summit.  Exact Lean names and modules are given in the companion
note and appendix.
\end{abstract}

\tableofcontents

% =============================================================================
\section{The problem}
\label{sec:problem}

\subsection{Outer certificates}

Consider any system that needs to certify, summarize, or represent complex
internal structure through an external or interface-level predicate.  This
is an extremely common structural situation.  Examples include:

\begin{itemize}
  \item A \emph{compressed witness} that claims to capture the essential
    content of a much richer proof or derivation.
  \item An \emph{abstraction layer} that claims to preserve enough structure
    for downstream reasoning to proceed faithfully.
  \item A \emph{quotient or interface representation} that projects internal
    distinctions into a coarser type and claims the projection is
    sufficient for the purposes at hand.
  \item A \emph{transport or portability claim}: a structure moved along a
    map or reindexing that claims to retain the relevant certificate content.
  \item An \emph{upgrade witness}: a certified object at one level of
    structure that claims to certify a richer level.
\end{itemize}

In all of these cases, the outer predicate plays the same role.  It claims
to stand in for richer internal content.  The question is when that claim
is honest.

\subsection{The adequacy burden}

The Adequacy Architecture formalizes a prior and more fundamental constraint
on this question.  Before asking \emph{which} outer certificates are honest,
one must understand what it means for any account to be adequate.

The core insight of the architecture is:

\begin{keybox}[title=The fundamental theorem (informal)]
  If a system is rich enough, adequacy has a cost.  You cannot freely
  flatten, summarize, certify, or route around the burden of adequacy
  without that burden reappearing somewhere else.
\end{keybox}

This is not a philosophical claim.  It is formalized as a hierarchy of
machine-checked theorems:

\begin{description}
  \item[\textbf{No free adequate collapse.}]
    If a representation is adequate in a non-trivial mode, it cannot erase
    all burden without paying elsewhere.  (\lean{no\_free\_adequate\_collapse},
    \leanmod{Burden/NoFreeCollapse.lean})
  \item[\textbf{Burden relocates; it does not vanish.}]
    Honest simplification moves structural obligations to another layer,
    carrier, or observability slot.  (\lean{burden\_not\_erased\_only\_relocated},
    \lean{adequacy\_preserving\_drop\_implies\_relocation},
    \leanmod{Burden/Relocation.lean})
  \item[\textbf{Route architecture is forced.}]
    Once adequacy and burden are genuinely entangled, route and finality
    claims stop being decorative --- they become part of what must be
    exhibited.  (\lean{master\_conditional\_summit\_ridge\_at},
    \leanmod{Lawful/MasterConditionalSummit.lean})
  \item[\textbf{Every map has a fundamental residual structure.}]
    Any forgetful or compressive morphism has a residual story: what was
    lost, where did it go, and what still observes it.
    (\lean{every\_map\_has\-\_fundamental\-\_residual\-\_package},
    \leanmod{Instances/WhatMapsForget.lean}; residual geometry via
    the residual-geometry library \lean{rcsOfMap} / \lean{vanishingCriterion})
  \item[\textbf{The highest reachable summit.}]
    $H_{\mathrm{final}}$: the highest reachable summit at any carrier is
    packaged as the theorem \lean{highestReachable\_summit\_at}
    (\leanmod{Lawful/FinalConditionalSummit.lean}).  This is the
    target that all outer admissibility routes aim to supply.  The
    three instantiation branches (halting/semantic, infinity-compression,
    and continuation) each supply $H_{\mathrm{final}}$ at the standard
    carrier; the outer gate unification theorem (\S\ref{sec:gate}) shows
    that all three routes into the gate are equivalent.
\end{description}

\subsection{The outer classification problem}

Given this foundation, the outer classification problem becomes precise.

In the Adequacy Architecture, the outer layer is the interface between the
abstract adequacy theory and concrete instantiations.  An outer certificate
predicate $C : \mathsf{Type} \to \mathsf{Prop}$ evaluated at a type-indexed
slot $\gamma$ claims: ``something meaningful is happening at this slot that
connects to the frontier of the adequacy landscape.''

The question is not merely ``does some predicate hold?''  It is:

\begin{enumerate}[label=(\arabic*)]
  \item \textbf{Which predicates are genuinely admissible?}
    Meaning: which $C$ imply the existence of an absolute frontier S1
    certificate?
  \item \textbf{Does admissibility unify?}
    Are the different presentations of admissibility actually the same gate?
  \item \textbf{What is the architecture above the gate?}
    Does admission unlock richer structure? Is that structure primitive or
    derived?
  \item \textbf{What are the honest limits?}
    Is there content that lies above both the gate and the second layer but
    is not part of either?
\end{enumerate}

The outer admissibility program of the Adequacy Architecture answers all four questions.

% =============================================================================
\section{Why this is hard and important}
\label{sec:importance}

\subsection{The na\"{\i}ve answer fails}

A na\"{\i}ve answer to the outer classification problem would be: any
predicate that happens to hold at some $\gamma$ counts as admissible.
That answer is empty --- it makes the notion of admissibility meaningless.

A better but still na\"{\i}ve answer would be: enumerate the known
certificate constructors and call each of them admissible.  But this
\emph{pluralist} answer raises its own problem: it produces a collection of
admissibility conditions without answering whether they form one class or
many, what their relationship is, and where the ceiling is.

\subsection{What makes the problem non-trivial}

The genuine difficulty is this.

From the supply side: the finite admissibility family $\Gadm$ is defined
by three structurally different constructor arms.  On the surface, these
look like three different kinds of certificates:
\begin{enumerate}
  \item A pullback display along a compare map.
  \item A full pack of conditions connecting the infinity-compression
    architecture to the standard carrier.
  \item A certified upgrade witness from the representation calculus.
\end{enumerate}
There is no obvious reason these should be the same.

From the demand side: any unified gate must have content.
It cannot be vacuous.  The witness $\Upb$ requires a genuine certified
frontier representation obtained by pulling back a lawful adequacy
architecture along a compare map.  This in turn requires a non-trivial
compression carrier pullback (through the infinity-compression arm) or
an involutive compare-map cleaving argument (through the carrier-swap arm),
each proved as a separate machine-checked theorem~\cite{SpivackAdequacy}.

\subsection{The breadth of the question}

The outer classification problem has breadth beyond the specific instances
examined.  It governs a wide family of certificate shapes:

\begin{itemize}
  \item Abstract pullback display certificates along compare maps.
  \item Finite admissibility presentations with multiple constructor arms.
  \item Grammar-generated outer classes under declared normal forms.
  \item Boundary certificates that are sound but fall outside the grammar.
  \item Functorial structural consequences of admission.
  \item Residual stronger conditions above both levels that are specific
    to a particular instantiation branch.
\end{itemize}

Each of these is a different kind of ``outer certificate.''  The theorem
stack answers where each kind fits in the architecture.

\subsection{The diagnostic value}

The five canonical laws of the Adequacy Architecture (informally: no free
collapse; burden relocates; route architecture is forced; adequate final
simplicity is never cheaply flat; if the burden is invisible, suspect the
map) all concern what happens at the interface between internal structure
and external claims.  The outer classification problem is a precise version
of this general question in the context of frontier certificates.

When someone presents an outer predicate as an admissibility condition, the
architecture asks:
\begin{itemize}
  \item Does it imply an absolute frontier S1 certificate?
  \item Is it grammar-generated, or does it sit outside the declared grammar?
  \item Does it unlock a second layer, or is it itself a second-layer
    certificate?
  \item Is it a primitive outer mode, or residual richer content?
\end{itemize}
Without answers to these questions, the admissibility claim is not fully
characterized.

% =============================================================================
\section{The conceptual thesis}
\label{sec:thesis}

Before the technical machinery, the central picture can be stated plainly.

\begin{keybox}[title=The main thesis (to be made precise in \S\S\ref{sec:gate}--\ref{sec:residual})]
\begin{enumerate}[label=(\arabic*)]
  \item There is a \textbf{single unified outer gate of honest admission}.
    Multiple finite admissibility presentations collapse to one gate
    $\Gadm \equiv \Upb$.
  \item The unified gate \textbf{has a real consequence}: it implies the
    existence of an absolute frontier S1 certificate.
  \item The gate is \textbf{maximal under the declared grammar}
    (NF0--NF3).  No grammar-generated class exceeds it.
  \item The \textbf{nearest expansion candidate} (NF4, the $\SoneLawful$
    predicate) does not force grammar expansion.  The boundary is stable.
  \item Above the gate, a \textbf{functorial second layer} is canonically
    produced without additional hypotheses.  It is not a new admissibility
    mode; it is what admission unlocks.
  \item \textbf{Stronger branch-specific content is residual}: particular
    to one instantiation branch, partially reconstructible from named
    hypotheses, and carrying a genuine irreducible obstruction.
    It is not primitive to the outer gate.
\end{enumerate}
\end{keybox}

The result is a three-level architecture:
\begin{itemize}
  \item \textbf{Level 1.} The outer gate of admission ($\Gadm \equiv \Upb$).
  \item \textbf{Level 2.} The functorial second layer ($\Upb \Rightarrow \Stwo$).
  \item \textbf{Residual band.} A branch-specific stronger condition above
    both levels --- partially reconstructible but carrying an irreducible
    obstruction not reachable from the outer gate alone.
\end{itemize}

This picture is not a sketch.  Every arrow is a named Lean theorem.  The
boundaries are not informal --- they are precisely characterized by what
does and does not follow from $\Upb$ alone.

% =============================================================================
\section{The formal world and disciplined scope}
\label{sec:scope}

\subsection{The type-indexed outer world}

Fix a universe of types.  An \emph{outer certificate predicate} is a
predicate $C : \mathsf{Type} \to \mathsf{Prop}$.  We evaluate such
predicates at a type-indexed slot $\gamma$.

\begin{definition}[Finite admissibility gate --- $\Gadm$]
  $\Gadm\,\gamma$ (\lean{Program1FiniteGAdm}) holds when $\gamma$ is
  witnessed by one of three constructor arms:
  \begin{enumerate}
    \item A pullback display arm:
      $\lean{Program1AdmissibilityPullbackDisplayWitness}\,\gamma$.
    \item A branch-native full pack arm (connecting the
      infinity-compression architecture to the standard carrier):
      $\lean{Program1IcStageDFullPackWitness}\,\gamma$.
    \item A representation-calculus upgrade witness arm:
      $\lean{Program1CertifiedUpgradeWitness}\,\gamma$.
  \end{enumerate}
\end{definition}

\begin{definition}[Pullback display witness --- $\Upb$]
  $\Upb\,\gamma$ (\lean{Program1Admissibility\-PullbackDisplayWitness}) holds
  when there exists a certified frontier representation obtained by
  pulling back a lawful adequacy architecture along a compare map at
  $\gamma$.
\end{definition}

$\Upb$ is not vacuous.  The branch-native arm requires a non-trivial
compression carrier pullback.  The carrier-swap arm requires an
involutive compare-map cleaving argument.  Both are proved as separate
machine-checked theorems~\cite{SpivackAdequacy}.

\subsection{The outer normal-form grammar}

The outer grammar is the closure of the following four normal forms under
declared grammar rules:
\begin{description}
  \item[NF0.] Bare pullback display.
  \item[NF1.] Certified frontier row.
  \item[NF2.] Upgrade witness.
  \item[NF3.] Composition of the above.
\end{description}
$\Ggen\,C$ means $C : \mathsf{Type} \to \mathsf{Prop}$ is
generated by this grammar.  NF4 ($\SoneLawful$), the lawful S1 frontier
row predicate, lies outside the grammar.  The maximality theorem
(\S\ref{sec:maximality}) shows that nothing grammar-generated exceeds
the outer gate, and the NF4 boundary analysis shows that even this
nearest non-generated candidate does not force expansion.

\subsection{Notation}

\begin{center}
\small
\begin{tabular}{lll}
\hline
Symbol & Lean name & Meaning \\
\hline
$\Gadm\,\gamma$ & \lean{Program1FiniteGAdm} & finite admissibility gate \\
$\Upb\,\gamma$ & \lean{Program1AdmissibilityPullbackDisplayWitness} & pullback display witness \\
$\Ggen\,C$ & \lean{Program1GrammarGeneratedOuterGate} & grammar-generated \\
$\AbsRawSone(P,B)$ & \lean{AbsoluteFrontierRawS1} & absolute frontier S1 certificate \\
$\SoneLawful\,\gamma$ & \lean{Nonempty (S1LawfulFrontierRow $\gamma$)} & NF4 win predicate \\
$\Stwo\,\gamma$ & \lean{pullbackDisplay\_with\_host\_summitFE3} & S5 / Level-2 package \\
$\StageD\,\gamma$ & \lean{icSpec031\_stageD\_resolution\_pack} & IC Stage-D four-conjunct pack \\
\hline
\end{tabular}
\end{center}

\subsection{Prior summit theorem}

The outer gate program builds on an earlier machine-checked summit theorem:
\lean{highestReachable\_summit\_at} (\leanmod{Lawful/FinalConditionalSummit.lean})
packages the highest reachable summit at any carrier as a structured
machine-checked object $H_{\mathrm{final}}$.  This is the target that all
outer certificate routes aim to supply --- admission at $\gamma$ means
supplying $H_{\mathrm{final}}$ through the pullback display witness.  The
unification theorem (\S\ref{sec:gate}) shows that all three constructor
arms of $\Gadm$ supply this target through $\Upb$, so the three routes
to the summit are logically equivalent.

\subsection{Residual geometry}

The Lean package imports a residual-geometry library that formalizes what
forgetful morphisms lose.  The canonical theorem
\lean{every\_map\_has\_fundamental\_residual\_package}
(\leanmod{Instances/WhatMapsForget.lean}) states that every forgetful
morphism has a fundamental residual structure: a machine-checked record
of what was lost in the map, what still observes it, and what the residual
geometry looks like.  This is the machine-checked form of the fifth
canonical law (``if the burden is invisible, suspect the map''), and it
provides the general framework in which the residual band of
\S\ref{sec:residual} is a specific instance.

\subsection{Scope discipline}

\begin{itemize}
  \item No claim of unconstrained $\forall P\,B$.  An unconditional
    frontier result for arbitrary predicate pairs is the next genuine
    research frontier; it is not part of the present theorems (see
    \S\ref{sec:beyond} and~\cite{SpivackAdequacy}).
  \item Maximality is grammar-cut: within the declared normal-form grammar
    (NF0--NF3), not over all possible predicates.
  \item The NF4 boundary analysis is a boundary report, not a completeness
    claim.
  \item The residual band (\S\ref{sec:residual}) is branch-specific.  Its
    irreducibility is an honest finding, not a deficiency.
  \item All theorems below are machine-checked; 0 \texttt{sorry}.
    Toolchain: Lean~4 v4.29.0-rc6, Mathlib v4.29.0-rc6~\cite{Mathlib}.
\end{itemize}

% =============================================================================
\section{Collapse to the outer gate}
\label{sec:gate}

\subsection{The unification theorem}

The first major result shows that the apparent plurality of admissibility
routes collapses.

\begin{theorem}[Unification]
\label{thm:unification}
  For all $\gamma$,\quad $\Gadm\,\gamma \iff \Upb\,\gamma$.
\end{theorem}

\begin{leanbox}
program1FiniteGAdm_iff_program1AdmissibilityPullbackDisplayWitness :
  forall gamma, G_adm gamma <-> U_pullback gamma
-- Module: Instances/Program1MetaUnificationMU1.lean
\end{leanbox}

This is nontrivial in both directions.

\emph{Three arms collapse (the non-trivial direction).}  Each of the three
constructor arms of $\Gadm$ must be shown to produce a $\Upb$ witness:
\begin{itemize}
  \item The pullback arm is immediate (first constructor).
  \item The branch-native full pack arm threads through a non-trivial
    compression carrier pullback (the infinity-compression compare-map
    representation, proved separately in the Lean library).
  \item The upgrade witness arm factors through the certified summit-row
    representation via the representation-calculus gate.
\end{itemize}

\emph{The gate is not empty (the honesty direction).}  $\Upb$ requires a
genuine lawful adequacy architecture pulled back along a compare map at
$\gamma$.  Four concrete witnesses have been constructed --- using the
infinity-compression pullback, an identity-level corpus instantiation, a
carrier-swap corpus instantiation, and a parity quotient on a three-element
type --- each requiring distinct non-trivial arguments~\cite{SpivackAdequacy}.

\subsection{The absolute frontier consequence}

\begin{theorem}[Absolute frontier S1 from admission]
\label{thm:frontier}
  For all $\gamma$,\quad $\Gadm\,\gamma \Rightarrow \exists P\,B,\,
  \AbsRawSone(P,B)$.
\end{theorem}

\begin{leanbox}
program1FiniteGAdm_implies_exists_absoluteFrontierRawS1 :
  forall gamma, G_adm gamma -> exists P B, AbsoluteFrontierRawS1 P B
-- Module: Instances/Program1MetaUnificationMU1.lean
\end{leanbox}

This is the outer payoff: admission at any $\gamma$ witnesses the
non-emptiness of the absolute frontier.  The absolute frontier
$\AbsRawSone(P,B)$ is the frontier of the adequacy landscape: it says
the highest reachable summit exists for some concrete predicate--burden pair,
as formalized in the machine-checked absolute-frontier program~\cite{SpivackAdequacy}.

\subsection{Why the collapse is significant}

Before Theorem~\ref{thm:unification}, one might have said: ``the adequacy
architecture has three kinds of outer certificates.''  After it, one can
say: ``there is one outer gate, and these are three ways to enter it.''
The compression is real.  The three arms were designed to capture different
structural routes to admissibility.  The theorem says those routes converge.

This is a classification result.  It says: for the disciplined
outer admissibility world, there is a unique admissibility condition
in the sense that all finite presentations of it are equivalent.

% =============================================================================
\section{Maximality and the grammar boundary}
\label{sec:maximality}

\subsection{Grammar-cut maximality}

\begin{theorem}[Grammar-cut maximality]
\label{thm:maximality}
  Let $C : \mathsf{Type} \to \mathsf{Prop}$.  If $\Ggen\,C$, then
  $\forall\,\gamma,\; C\,\gamma \Rightarrow \Upb\,\gamma$.
\end{theorem}

\begin{leanbox}
program1_outer_grammar_cut_maxA :
  forall (C : Type -> Prop),
    Program1GrammarGeneratedOuterGate C ->
    forall gamma, C gamma -> U_pullback gamma
-- Module: Instances/Program1OuterGateGrammar.lean
\end{leanbox}

Theorem~\ref{thm:maximality} establishes that $\Upb$ is the \emph{ceiling}
of the declared grammar.  Any outer predicate generated by composing,
lifting, or grammar-producing from NF0--NF3 stays within $\Upb$.

Without this theorem, the unified gate would be a lower bound: something
that everything we know satisfies, but which could in principle be exceeded
by some more expressive grammar-generated class we have not tried.  With it,
the gate is tight: within the declared scope, there is no room above it.

\subsection{The NF4 boundary finding}

The nearest principled candidate for a broader outer class is the NF4
predicate $\SoneLawful\,\gamma$, which asserts the non-emptiness of a
lawful S1 frontier row at $\gamma$.

\begin{theorem}[NF4 outer soundness]
\label{thm:nf4-sound}
  For all $\gamma$,\quad $\SoneLawful\,\gamma \Rightarrow
  \lean{program1OuterLayerBLaw}\,\gamma$.
\end{theorem}

\begin{leanbox}
Program1OuterCertificateSound_program1OuterWinS1LawfulNonempty :
  forall gamma, S1Lawful gamma -> program1OuterLayerBLaw gamma
-- Module: Instances/Program1MaximalityMA1.lean
\end{leanbox}

The NF4 boundary analysis records three facts jointly:
\begin{enumerate}
  \item $\SoneLawful$ is outer-sound: it implies the outer layer law
    (Theorem~\ref{thm:nf4-sound}).
  \item $\SoneLawful$ is \emph{not} grammar-generated (not NF0--NF3).
  \item There is \emph{no proved universal} $\SoneLawful\,\gamma \Rightarrow \Upb\,\gamma$.
\end{enumerate}

NF4 therefore remains a logical pocket: sound where instantiated, but not
promoted to the operative outer gate.

\subsection{Counterexample: S1-shape is necessary}

Before stating what the NF4 boundary means, we note that an explicit
machine-checked counterexample shows the S1-shape condition is not vacuous:

\begin{leanbox}
not_s1LawfulFrontierRow_counterexample_pair :
  forall ..., Not (Nonempty (S1LawfulFrontierRow counterexamplePair))
-- Module: Portability/CertifiedFrontierRow.lean (S1ShapeNecessityCounterexample)
\end{leanbox}

There exist explicit pairs $(P, B)$ where no lawful S1 frontier row exists.
This means the S1-shape requirement in the architecture is load-bearing:
it is not trivially satisfied.  The grammar boundary separates something
real from something vacuous.

\subsection{What the NF4 boundary means}

The NF4 finding matters in two directions simultaneously.

\emph{Downward:} it confirms the architecture is not trivially
all-encompassing.  The grammar boundary is real.  There exist predicates
(like $\SoneLawful$) that are not grammar-generated and not universally
equal to $\Upb$.

\emph{Upward:} NF4 did not force expansion.  The nearest candidate for
``maybe there is a broader class'' turned out not to yield one.  The
architecture is stable at its current boundary --- not because we could not
find anything stronger, but because the nearest candidate was examined and
did not fit.

Together, Theorems~\ref{thm:maximality} and~\ref{thm:nf4-sound} give a
rigorous sense in which $\Upb$ is the right outer gate at the current level
of structure: maximal in the declared grammar, and stable against the
nearest boundary test.

% =============================================================================
\section{The second layer above admission}
\label{sec:second-layer}

\subsection{The second-layer consequence}

One of the most structurally important results is not about the outer gate
itself.  It is about what the outer gate unlocks.

\begin{theorem}[Functorial second layer]
\label{thm:second-layer}
  For all $\gamma$,\quad $\Upb\,\gamma \;\Rightarrow\; \Stwo\,\gamma$
  (and therefore $\Gadm\,\gamma \;\Rightarrow\; \Stwo\,\gamma$).
\end{theorem}

\begin{leanbox}
program1_layer2PullbackDisplay_with_host_summitFE3_of_U_pullback :
  forall gamma, U_pullback gamma -> L2 gamma
program1_layer2PullbackDisplay_with_host_summitFE3_of_program1FiniteGAdm :
  forall gamma, G_adm gamma -> L2 gamma
-- Module: Instances/Program1LayeredStructureLS1.lean
\end{leanbox}

The $\Stwo$ package (\lean{pullbackDisplay\_with\_host\_summitFE3}) combines
the pullback display with a host-side summit certificate.  It is
obtained by projecting from the same witness that gives $\Upb\,\gamma$
--- no additional hypotheses are required.

\subsection{Why this is structurally important}

The key structural point:

\begin{itemize}
  \item $\Stwo$ is \emph{not} a new constructor of $\Gadm$.  It is not a
    new way to be admitted.
  \item $\Stwo$ is \emph{not} grammar-generated (not NF0--NF4).
  \item $\Stwo$ does not imply $\Gadm$ (it is downstream).
  \item $\Stwo$ is canonically produced from any $\Upb$ witness, without
    choice.
\end{itemize}

The outer certificate world therefore has a strict vertical structure:

\begin{center}
\textbf{Level 1} (admission) lies below \textbf{Level 2} (functorial
consequence).  These are not the same level.
\end{center}

This matters because it shows that the outer gate is not a dead end.
Admission at $\gamma$ is the beginning of a richer theory, not just a
yes/no question.  The architecture has vertical structure: gate below,
functorial consequence above.

\subsection{What it means to have a second layer}

Before this result, one might have asked: ``once we know a certificate is
admissible, is there more?''  The answer is yes, and specifically: there is
a canonical functorial package $\Stwo$ that comes automatically with any
admission witness.

This has a conceptual implication for how one should think about outer
certificates.  An outer certificate is not just a stamp of approval.  It
is an entry ticket to a richer structural world.  The second-layer package
packages the pullback display together with a host-side summit certificate,
giving a richer simultaneous view than admission alone provides.

% =============================================================================
\section{Stage D and the residual band}
\label{sec:residual}

\subsection{What the residual condition is}

The residual condition $\StageD$ is a branch-specific pack of requirements
connecting the infinity-compression architecture to the highest reachable
summit at the standard carrier.  It decomposes into four conjuncts:

\begin{theorem}[Residual condition decomposition]
\label{thm:staged-d}
  For all $\gamma$,
  \[
    \StageD\,\gamma
    \;\iff\;
    F_1(\gamma) \;\wedge\; F_2(\gamma) \;\wedge\; F_3(\gamma) \;\wedge\; F_4(\gamma)
  \]
  where $F_1 =$ obstruction entry, $F_2 =$ absolute-frontier S1
  conjunct, $F_3 =$ bypass, $F_4 =$ cascade.
\end{theorem}

Each of the four conjuncts has a named machine-checked lemma (see the
companion note for exact names).  Facets $F_1$, $F_3$, $F_4$ are
reconstructible from named hypotheses.  $F_2$ is the irreducible residue.

\subsection{The residuality result}

\begin{proposition}[Residual condition residuality]
\label{prop:sr-b}
  \begin{enumerate}[label=(\alph*)]
    \item $\Upb\,\gamma$ does not imply $\StageD\,\gamma$ in general.
    \item $\StageD\,\gamma$ does not imply $\Upb\,\gamma$ without the $F_1$
      and bypass arms.
    \item $F_1$, $F_3$, $F_4$ are reconstructible from named hypotheses.
    \item $F_2$ (the absolute-frontier S1 conjunct) is the irreducible
      residue: it does not follow from $\Upb\,\gamma$ alone.
  \end{enumerate}
\end{proposition}

\subsection{Why residuality matters}

Without this residuality result, one might conclude: ``the residual
condition is just another outer admissibility condition; we should add it
to the gate.''  Proposition~\ref{prop:sr-b} closes this off decisively.

The residual condition $\StageD$ is not primitive to the outer gate.  It
is not part of the second layer.  It is branch-specific extra structure
that has its own internal organization (the four conjuncts) and its own
obstruction ($F_2$).

This means the architecture has a clean outer boundary ($\Upb$) and a
clean second layer ($\Stwo$), and \emph{everything beyond that is
residual} --- not foreign, not ignored, but structurally downstream
and branch-specific.

The residuality result is what makes the two-level picture honest: it says
that the genuinely stronger branch-specific conditions do not sneak back
into Level~1 or Level~2.  They live in their own residual band, partially
accessible via named lemmas, with an irreducible core.

\subsection{The residual band and what maps forget}

The residual band is not an ad hoc finding.  It is the outer gate
manifestation of the general residual principle formalized throughout
the architecture.

The theorem \lean{every\_map\_has\_fundamental\_residual\_package}
states that every forgetful morphism has a fundamental residual
structure --- a machine-checked record of what was lost in the map, what
still observes it, and what the residual geometry looks like.  This is
the machine-checked form of the fifth canonical law (``if the burden is
invisible, suspect the map'').

The irreducible conjunct $F_2$ of the residual condition is a specific
instance of this principle: the map from the outer gate to the
infinity-compression architecture forgets exactly the $F_2$ content, and
that content is what the fundamental residual structure records.

\subsection{The significance of the residual band}

The existence of a residual band is not a failure of the architecture.
It is a finding.

The architecture is not claiming that $\Upb$ governs everything.  It is
claiming that $\Upb$ is the right outer gate at its level of structure,
and that what lies beyond it is branch-specific content with its own
structural character.  The residuality result locates that content precisely.

This is the deeper picture.  The residual band is the boundary between
``what the outer gate governs'' and ``what requires more specific internal
structure to access.''  Knowing where that boundary is, and why it is
there, is part of the architecture's contribution.

% =============================================================================
\section{The complete architecture}
\label{sec:architecture}

\subsection{The three-level picture}

\begin{center}
\begin{tcolorbox}[colback=gray!5, colframe=gray!40, width=0.88\linewidth,
  arc=3pt, boxrule=0.5pt]
\begin{Verbatim}[fontsize=\small]
  NF4 (S1Lawful)  -- outer-sound; not grammar-generated;
                     no universal S1Lawful -> U_pullback
        ^
  ---------------------------------------------------------
  Residual: branch-specific condition  (partial basis;
                       F2 is irreducible w.r.t. U_pullback)
        ^
  ---------------------------------------------------------
  Level 2: L2 = pullbackDisplay_with_host_summitFE3
                [functorial from U_pullback; not a new outer mode]
        ^
  ---------------------------------------------------------
  Level 1: G_adm <=> U_pullback  (outer gate; maximal NF0--NF3)
        v
  exists P B, AbsoluteFrontierRawS1 P B
\end{Verbatim}
\end{tcolorbox}
\end{center}

Every arrow is a named machine-checked theorem.  The horizontal rules
are logical boundaries characterized by what is and is not derivable
from $\Upb$ alone.

\subsection{The integrated claim}

For the disciplined outer admissibility world:

\begin{enumerate}[label=(\arabic*)]
  \item \textbf{One outer gate.}
    $\Gadm \equiv \Upb$: the three structurally different admissibility
    arms collapse to the same gate
    (Theorem~\ref{thm:unification}).

  \item \textbf{Real consequence.}
    The gate implies the existence of an absolute frontier S1 certificate
    (Theorem~\ref{thm:frontier}).

  \item \textbf{Grammar ceiling.}
    No grammar-generated predicate (NF0--NF3) exceeds $\Upb$
    (Theorem~\ref{thm:maximality}).

  \item \textbf{Stable boundary.}
    The nearest expansion candidate (NF4) does not force grammar expansion;
    the architecture is stable at its current scope
    (\S\ref{sec:maximality}).

  \item \textbf{Functorial second layer.}
    $\Upb\,\gamma \Rightarrow \Stwo\,\gamma$, canonically, from the same
    witness, without additional hypotheses
    (Theorem~\ref{thm:second-layer}).

  \item \textbf{Residual band.}
    $\StageD$ is not primitive to the outer gate, carries a genuine
    branch-specific irreducible residue, and is partially reconstructible
    from named lemmas
    (Proposition~\ref{prop:sr-b}).
\end{enumerate}

\subsection{What every arrow earns}

The diagram above is not a schematic.  Each connection has a specific
machine-checked proof:

\begin{description}
\item[] $\Gadm \iff \Upb$:
    \lean{program1FiniteGAdm\_iff\_program1Admissibility\-PullbackDisplayWitness}
  \item[] $\Gadm \Rightarrow \AbsRawSone$:
    \lean{program1FiniteGAdm\_implies\_exists\_absolute\-FrontierRawS1}
  \item[] $\Ggen \subseteq \Upb$:
    \lean{program1\_outer\_grammar\_cut\_maxA}
  \item[] $\SoneLawful \Rightarrow \text{outer soundness}$:
    \lean{Program1OuterCertificate\-Sound\_program1Outer\-WinS1LawfulNonempty}
  \item[] $\Upb \Rightarrow \Stwo$:
    \lean{program1\_layer2PullbackDisplay\_with\_host\_\-summitFE3\_of\_U\_pullback}
  \item[] $\StageD$ residuality:
    \lean{icSpec031\_stageD\_resolution\_pack} + facet lemmas
\end{description}

% =============================================================================
\section{Robustness and breadth of applicability}
\label{sec:robustness}

\subsection{Multiple arms already collapsed}

The unification result covers a concrete range of admissibility
presentations already:

\begin{itemize}
  \item The infinity-compression pullback arm (compression carrier along
    a constant compare map).
  \item Two carrier-level corpus instantiations (identity and carrier-swap,
    the latter proved via an involutive cleaving argument).
  \item A parity quotient on a three-element type (a strict intermediate
    quotient on $\mathrm{Fin}\,3$).
  \item The certified upgrade witness arm (representation-calculus gate).
\end{itemize}

These are structurally heterogeneous instances.  That they all collapse to
$\Upb$ is not obvious from their definitions; each required a distinct
proof.

\subsection{Why the maximality result is not local}

Grammar-cut maximality says: \emph{any} predicate generated by the outer
normal-form grammar collapses to $\Upb$.  The grammar is not just about
the specific instances above.  It is about the closure under NF0--NF3.
This means the result generalizes beyond the specific instantiations
proved so far, to all possible grammar-generated predicates.

\subsection{The boundary test}

The NF4 boundary finding provides an independent robustness check.  Before
claiming that $\Upb$ is the right outer gate, one should ask: is there
some nearby candidate that would exceed it?  The nearest candidate was
$\SoneLawful$ (the NF4 win predicate).  That candidate was examined and
did not force expansion.  This provides evidence that the grammar boundary
is genuine rather than accidental.

\subsection{Honest scope}

The architecture does not claim to govern all possible outer predicates.
That claim would be false: NF4 exists and is not grammar-generated.  The
correct claim is grammar-cut maximality, which is both true and
informative.

Similarly, the architecture does not claim that the residual condition is
unreachable or unimportant.  The residuality result says $\StageD$ is
residual --- it has a partial basis and an irreducible residue.  That is
an honest characterization of its relationship to the outer gate, not a
dismissal.

% =============================================================================
\section{Breadth: what this architecture governs}
\label{sec:breadth}

\subsection{The class of systems it applies to}

The Adequacy Architecture governs any system where:
\begin{itemize}
  \item There is a notion of adequacy with real, non-erasable burden.
  \item There is a lawful class of architectures supporting stratified
    adequacy, bottleneck structure, and bridge relationships between layers.
  \item There are outer certificate predicates evaluated at type-indexed
    slots.
  \item There is a grammar of outer normal forms.
\end{itemize}

Within such a system, the entire theorem stack applies.

\subsection{Three instantiation branches}

The architecture has three concrete instantiation branches:
\begin{itemize}
  \item The \emph{halting/semantic} branch (NEMS), which provides
    halting-anchored semantic summit certificates.
  \item The \emph{infinity-compression} branch (IC), which provides
    compression carriers and native stage-D structure.
  \item The \emph{continuation uniformization} branch (APS), which
    provides indexed middle-layer structure via recursion uniformization.
\end{itemize}

Each branch supplies a distinct instantiation route into the outer gate.
The three-branch closure theorem
(\lean{three\_branch\_corpus\_highest\_summit\_at})
shows that all three branches supply the highest reachable summit at the
standard carrier.  The outer gate unification (Theorem~\ref{thm:unification})
then shows that what all three branches contribute collapses to the same
outer gate.

This means the outer gate theorem is not a statement about one branch.
It is a statement about the whole three-branch structure.

The halting/semantic branch carries a full semantic faithfulness story.
The theorem \lean{summitReportOf\-\_indexed\-\_master\-\_agreement\-\_consequences}
(\leanmod{Instances/NEMSSemantic\-FaithfulTransport.lean})
packages all consequences of master-indexed summit agreement --- agreed
summit, summit certificate, indexed certificates, forgetful lifts through
the fourfold summit slot --- as a single machine-checked bundle.
This is the branch-level analogue of the outer gate: the halting/semantic
faithfulness story collapses to a single master agreement consequence,
just as the outer gate collapses to a single $\Upb$.

\subsection{The five canonical laws, re-read}

The five canonical diagnostic laws of the Adequacy Architecture (no free
collapse; burden relocates; route architecture is forced; adequate final
simplicity is never cheaply flat; if the burden is invisible, suspect
the map) each find a reflection in the outer classification results:

\begin{enumerate}
  \item \emph{No free collapse}: the outer gate is not vacuous and cannot
    be entered without a genuine pullback witness.
  \item \emph{Burden relocates}: the three admissibility arms are not
    equivalent by default --- the pullback structure is what they all
    relocate their structure into.
  \item \emph{Route architecture is forced}: the second layer $\Stwo$ is
    the route architecture that admission forces.
  \item \emph{Final simplicity is not cheaply flat}: NF4 exists, is sound,
    but does not collapse to the gate.  The nearest ``simpler'' account
    is not equivalent to the gate.
  \item \emph{If the burden is invisible, suspect the map}: the residual
    condition's irreducibility ($F_2$) is precisely this --- the map from
    the outer gate to the branch-specific architecture forgets exactly
    the $F_2$ content, which is where the residual burden is held.
\end{enumerate}

% =============================================================================
\section{What lies beyond the summit}
\label{sec:beyond}

The architecture is not claiming to be complete in a global sense.  Three
frontiers are honestly open.

\subsection{Global unconditional frontier}

The existence of an absolute frontier S1 certificate for an arbitrary
predicate--burden pair without hypotheses --- $\AbsRawSone(P,B)$ for all
$P,B$ --- remains open.  This is the next genuine research frontier
beyond the outer gate program~\cite{SpivackAdequacy}.  It is not a missing
piece of the current architecture; it is where the current architecture
points next.

\subsection{Broader outer grammar}

If a principled grammar larger than the current NF0--NF3 normal-form
grammar were ever chartered, the grammar-cut maximality result would need
re-examination in that scope.  The NF4 boundary finding is not a proof that
no broader grammar is possible.  It is a proof that the nearest specific
candidate (NF4) did not force one now.  A future investigation could, in
principle, extend the grammar and re-examine the ceiling.

\subsection{Residual condition universality}

Whether the residual condition $\StageD$ could ever be brought inside the
outer gate for all $\gamma$ --- i.e., whether $\Upb \Rightarrow \StageD$
universally --- or whether the branch-specific obstruction ($F_2$) is
genuinely permanent, is open.  The residuality result gives ``partial
basis, irreducible residue.''  Establishing $\Upb \Rightarrow \StageD$
universally would be a genuine extension of the current theorem stack.

\subsection{What these open frontiers mean}

These three frontiers are not gaps in the current theorem stack.  They are
the next genuine research questions, now clearly located by the architecture.
The architecture does not hide them; it exposes them precisely.

A frontier that is clearly located is more valuable than a frontier that is
obscured by vague language.  The current theorem stack has the virtue of
telling you exactly where the edge is and why.

% =============================================================================
\section{Machine verification and reproducibility}
\label{sec:machine}

\subsection{The machine-checked status}

All theorems stated in this paper are machine-checked in Lean 4 with 0
\lean{sorry}.  The toolchain is \lean{leanprover/lean4:v4.29.0-rc6} with
Mathlib \lean{v4.29.0-rc6}.

To reproduce: run \lean{lake build} from
\lean{adequacy-architecture-lean/}.  The build succeeds with 0 sorry.
The reproducibility script is \lean{scripts/verify-lean-build.sh}.

\subsection{The companion note}

The companion note (\emph{Technical Companion: The Solved Theorem Stack of
the Adequacy Architecture}) gives the same theorems in a concise,
theorem-dense format suitable for technically fluent readers who need exact
Lean names and proof-sketch references without the motivating framing of
this paper.

The two documents are designed to be read together.  This paper provides
altitude and conceptual orientation; the companion provides precision and
artifact accountability.

% =============================================================================
\section*{Acknowledgements}

All machine-checked proofs are in the \texttt{adequacy-architecture-lean}
repository.  A full theorem-name index and reproducibility instructions
are included in that repository.  See~\cite{SpivackAdequacy} for the
complete program record and the mathematical kernel statement.

% =============================================================================
\appendix

\section{Theorem table: outer gate stack}
\label{sec:thmtable}

All modules are under
\texttt{adequacy-architecture-lean/AdequacyArchitecture/}.
All theorems machine-checked; 0 \texttt{sorry}.
The \textbf{SPEC} column gives the internal development reference
(for cross-referencing the repository).

\begingroup\footnotesize
\setlength{\tabcolsep}{3pt}
\renewcommand{\arraystretch}{1.15}
\begin{tabular}{@{}lp{0.50\linewidth}p{0.28\linewidth}r@{}}
\hline
\textbf{Result} & \textbf{Lean name} & \textbf{Module} & \textbf{SPEC} \\
\hline
Thm.~\ref{thm:unification} &
  \raggedright\texttt{program1FiniteGAdm\_iff\_program1Admissibility\-PullbackDisplayWitness} &
  \raggedright\texttt{Program1MetaUnification\-MU1.lean} & 037 \\
Thm.~\ref{thm:frontier} &
  \raggedright\texttt{program1FiniteGAdm\_implies\_exists\_absolute\-FrontierRawS1} &
  \raggedright\texttt{Program1MetaUnification\-MU1.lean} & 036 \\
Thm.~\ref{thm:maximality} &
  \raggedright\texttt{program1\_outer\_grammar\_cut\_maxA} &
  \raggedright\texttt{Program1OuterGate\-Grammar.lean} & 041 \\
Thm.~\ref{thm:nf4-sound} &
  \raggedright\texttt{Program1OuterCertificateSound\_\-program1OuterWinS1LawfulNonempty} &
  \raggedright\texttt{Program1Maximality\-MA1.lean} & 044 \\
Thm.~\ref{thm:second-layer} &
  \raggedright\texttt{program1\_layer2PullbackDisplay\_with\_host\_\-summitFE3\_of\_U\_pullback} &
  \raggedright\texttt{Program1Layered\-StructureLS1.lean} & 042 \\
Thm.~\ref{thm:second-layer} &
  \raggedright\texttt{program1\_layer2PullbackDisplay\_with\_host\_\-summitFE3\_of\_program1FiniteGAdm} &
  \raggedright\texttt{Program1Layered\-StructureLS1.lean} & 042 \\
Thm.~\ref{thm:staged-d} &
  \raggedright\texttt{icSpec031\_stageD\_resolution\_pack} &
  \raggedright\texttt{ProgramFiniteAdmissi\-bility.lean} & 031 \\
Prop.~\ref{prop:sr-b} $F_1$ &
  \raggedright\texttt{icSpec031\_stageD\_iff\_conjunct\_of\_af1\_\-native\_carrier} &
  \raggedright\texttt{ProgramFiniteAdmissi\-bility.lean} & 043 \\
Prop.~\ref{prop:sr-b} $F_3$ &
  \raggedright\texttt{icSpec031\_stageD\_bypass\_conjunct\_of\_af1\_\-standard} &
  \raggedright\texttt{ProgramFiniteAdmissi\-bility.lean} & 043 \\
Prop.~\ref{prop:sr-b} $F_4$ &
  \raggedright\texttt{icSpec031\_stageD\_cascade\_conjunct\_of\_af1\_\-standard} &
  \raggedright\texttt{ProgramFiniteAdmissi\-bility.lean} & 043 \\
\S\ref{sec:maximality} &
  \raggedright\texttt{not\_s1LawfulFrontierRow\_counterexample\_pair} &
  \raggedright\texttt{Portability/Certified\-FrontierRow.lean} & 030 \\
\hline
\end{tabular}
\endgroup

\bigskip

\section{Theorem table: foundational theorems}
\label{sec:landmarks}

Selected foundational theorems that the outer gate program builds on.
All machine-checked; 0 \texttt{sorry}.

\begingroup\footnotesize
\setlength{\tabcolsep}{3pt}
\renewcommand{\arraystretch}{1.15}
\begin{tabular}{@{}p{0.46\linewidth}p{0.38\linewidth}r@{}}
\hline
\textbf{Lean name} & \textbf{Module} & \textbf{SPEC} \\
\hline
\raggedright\texttt{adequacy\_interface\_exists} &
  \raggedright\texttt{Foundation/Adequacy\-Interface.lean} & A1 \\
\raggedright\texttt{carrier\_exists\_of\_adequacy\_mode} &
  \raggedright\texttt{Foundation/BurdenCarrier\-Existence.lean} & A2 \\
\raggedright\texttt{internality\_outsourcing\_barrier} &
  \raggedright\texttt{Foundation/Internality\-Outsourcing.lean} & A3 \\
\raggedright\texttt{no\_free\_adequate\_collapse} &
  \raggedright\texttt{Burden/NoFreeCollapse.lean} & B1 \\
\raggedright\texttt{burden\_not\_erased\_only\_relocated} &
  \raggedright\texttt{Burden/Relocation.lean} & B2 \\
\raggedright\texttt{adequacy\_preserving\_drop\_implies\_relocation} &
  \raggedright\texttt{Burden/Relocation.lean} & B3 \\
\raggedright\texttt{universal\_irreducible\_adequacy\_lawful} &
  \raggedright\texttt{Lawful/MasterTheorem.lean} & 008 \\
\raggedright\texttt{bottleneck\_one\_bundle} &
  \raggedright\texttt{Lawful/Bottleneck\-Theorems.lean} & 009 \\
\raggedright\texttt{master\_conditional\_summit\_ridge\_at} &
  \raggedright\texttt{Lawful/MasterConditional\-Summit.lean} & 021 \\
\raggedright\texttt{highestReachable\_summit\_at} &
  \raggedright\texttt{Lawful/FinalConditional\-Summit.lean} & 022 \\
\raggedright\texttt{three\_branch\_corpus\_highest\_summit\_at} &
  \raggedright\texttt{Instances/ThreeBranch\-Closure.lean} & 026 \\
\raggedright\texttt{standardRoleAwareRepresentationBundle} &
  \raggedright\texttt{Instances/RoleAwareThree\-BranchFrontier.lean} & 030 \\
\raggedright\texttt{AbsoluteFrontierRawS1}, \texttt{HFrontier} &
  \raggedright\texttt{AbsoluteFrontier/} & 027 \\
\raggedright\texttt{every\_map\_has\_fundamental\_residual\_package} &
  \raggedright\texttt{Instances/WhatMaps\-Forget.lean} & (imported) \\
\raggedright\texttt{lawfulAdequacyArchitecture\_\-pullbackAlongCompare} &
  \raggedright\texttt{Lawful/ComparePullback.lean} & 031 \\
\raggedright\texttt{icSpec031\_stageD\_resolution\_pack} (full four-conjunct decomp.) &
  \raggedright\texttt{Instances/ICCompareRepresentation\-Pullback.lean} & 031 \\
\raggedright\texttt{CertifiedSummitRowRepresentation} &
  \raggedright\texttt{Portability/Representation\-CalculusRL9.lean} & 033 \\
\raggedright\texttt{certifiedSummitRowRepr\_summitFE3\_\-joint\_upgrade\_iff\_row} &
  \raggedright\texttt{Portability/Representation\-CalculusRL9.lean} & 033 \\
\raggedright\texttt{summitReportOf\_indexed\_master\_\-agreement\_consequences} &
  \raggedright\texttt{Instances/NEMSSemantic\-FaithfulTransport.lean} & 025 \\
\raggedright\texttt{M1\_from\_certifiedFrontierRow} &
  \raggedright\texttt{Portability/Spec034Program\-ClosureLaw.lean} & 034 \\
\raggedright\texttt{fe4\_earned\_absoluteFrontierRawS1\_ladder\_bundle} &
  \raggedright\texttt{Portability/AbsoluteFrontier\-EarnedToy.lean} & 029 \\
\raggedright\texttt{program1FiniteGAdm\_icNems} &
  \raggedright\texttt{Instances/ProgramFinite\-Admissibility.lean} & 036 \\
\hline
\end{tabular}
\endgroup

% =============================================================================
\section*{References}

\begin{thebibliography}{99}

\bibitem{SpivackAdequacy}
N.~Spivack.
\newblock \emph{Adequacy Architecture: Outer Gate, Layered Certificate
  Architecture, and Machine-Checked Synthesis}.
\newblock Adequacy Architecture Project, \texttt{adequacy-architecture-lean}
  repository, 2026.

\bibitem{SpivackIC}
N.~Spivack.
\newblock \emph{Canonical Certification Does Not Exhaust Reflective Structure:
  Reflective Split, Strict Refinement, and Fiber Structure in Infinity
  Compression}.
\newblock Flagship manuscript, \texttt{infinity-compression} repository, 2026.

\bibitem{SpivackSummit}
N.~Spivack.
\newblock \emph{Reflective Non-Exhaustion: Certification, Realization, and the
  Summit Program}.
\newblock Summit manuscript, \texttt{infinity-compression} repository, 2026.

\bibitem{SpivackSynthesis}
N.~Spivack.
\newblock \emph{Certification, Realization, and Obstruction: A Universal Fiber
  Architecture}.
\newblock Synthesis manuscript, \texttt{infinity-compression} repository, 2026.

\bibitem{SpivackValidation}
N.~Spivack.
\newblock \emph{External Validation of a Positive-Closure Proof Architecture:
  A Multi-Tranche Survey}.
\newblock Validation manuscript, \texttt{infinity-compression} repository, 2026.

\bibitem{SpivackAPS}
N.~Spivack.
\newblock \emph{APS Recursion Uniformization: Composition, Gluing, and
  Separation of Fixed-Point Theorems}.
\newblock \texttt{aps-recursion-uniformization-lean} repository, 2024--2026.

\bibitem{SpivackFiber}
N.~Spivack.
\newblock \emph{Fiber Architecture for Group Extensions in Lean~4: Cocycles,
  Splitting, and the Cohomological Bridge}.
\newblock General-method manuscript, \texttt{infinity-compression} repository, 2026.

\bibitem{SpivackQuillen}
N.~Spivack.
\newblock \emph{Quillen's Theorem~A for Galois Connections: A Machine-Checked
  Formalization in Lean~4}.
\newblock General-method manuscript, \texttt{infinity-compression} repository, 2026.

\bibitem{Lean4}
L.~de~Moura and S.~Ullrich.
\newblock The {L}ean~4 theorem prover and programming language.
\newblock In A.~Platzer and G.~Sutcliffe, editors,
  \emph{Automated Deduction --- CADE~28}, volume 12699 of \emph{LNCS},
  pages 625--635. Springer, 2021.

\bibitem{Mathlib}
The Mathlib community.
\newblock The {L}ean mathematical library.
\newblock In \emph{Proceedings of the 9th ACM SIGPLAN International Conference
  on Certified Programs and Proofs (CPP~2020)}, pages 367--381. ACM, 2020.
\newblock \url{https://github.com/leanprover-community/mathlib4}.

\bibitem{HoTTBook}
The Univalent Foundations Program.
\newblock \emph{Homotopy Type Theory: Univalent Foundations of Mathematics}.
\newblock Institute for Advanced Study, 2013.

\bibitem{Rice1953}
H.~G.~Rice.
\newblock Classes of recursively enumerable sets and their decision problems.
\newblock \emph{Transactions of the American Mathematical Society},
  74(2):358--366, 1953.

\bibitem{Turing1936}
A.~M.~Turing.
\newblock On computable numbers, with an application to the
  Entscheidungsproblem.
\newblock \emph{Proceedings of the London Mathematical Society} (2)
  42:230--265, 1936.

\bibitem{MacLane}
S.~Mac~Lane.
\newblock \emph{Categories for the Working Mathematician}.
\newblock Springer, 2nd edition, 1998.

\bibitem{Hatcher2002}
A.~Hatcher.
\newblock \emph{Algebraic Topology}.
\newblock Cambridge University Press, 2002.

\end{thebibliography}

\end{document}
